% chap10_matplotlib.tex - Data visualization with Matplotlib
% ============================================================================

\chapter{Matplotlib: Visualization}
\label{chap:matplotlib}

Clear figures are essential for understanding data and communicating results. Matplotlib is Python's foundational plotting library---flexible, powerful, and capable of producing publication-quality graphics.

% ----------------------------------------------------------------------------
\section{The Basics}
\label{sec:plotting-basics}

\begin{lstlisting}
import matplotlib.pyplot as plt
import numpy as np

# Create data
x = np.linspace(0, 10, 100)
y = np.sin(x)

# Create plot
plt.plot(x, y)
plt.show()
\end{lstlisting}

\begin{center}
\includegraphics[width=0.7\textwidth]{figures/chap10/basic_sine.pdf}
\end{center}

This creates a simple line plot. \py{plt.show()} displays the figure in a window or inline (in Jupyter notebooks). In spyder there is a "plot" pane which will keep a history of that session's plots for perusing.

\subsection{Adding Labels and Title}

\begin{lstlisting}
plt.plot(x, y)
plt.xlabel("Time (s)")
plt.ylabel("Amplitude")
plt.title("Sine Wave")
plt.show()
\end{lstlisting}

\begin{center}
\includegraphics[width=0.7\textwidth]{figures/chap10/sine_labeled.pdf}
\end{center}

\subsection{Saving Figures}

\begin{lstlisting}
plt.plot(x, y)
plt.xlabel("Time (s)")
plt.ylabel("Amplitude")
plt.savefig("sine_wave.png", dpi=300)  # 300 DPI for print quality
plt.savefig("sine_wave.pdf")            # vector format for publications
\end{lstlisting}

\begin{note}
\py{plt.savefig()} saves the figure in the same folder as the script. In the window opened by plt.show(), it is possible to choose another destination folder as well as the format and the name.
\end{note}

\begin{tip}
Save as PDF or SVG for publications---these are vector formats that scale without pixelation. For more information, read up on the difference between raster and vector graphics.
\end{tip}

\begin{warning}
Saving a figure overwrites the one with the same name in that folder (if it exists). Be careful when choosing how to name your figure files and especially if you copy-paste code within the same script.
\end{warning}

% ----------------------------------------------------------------------------
\section{Line Plots}
\label{sec:line-plots}

\subsection{Multiple Lines}

\begin{lstlisting}
x = np.linspace(0, 10, 100)

plt.plot(x, np.sin(x), label="sin(x)")
plt.plot(x, np.cos(x), label="cos(x)")
plt.xlabel("x")
plt.ylabel("y")
plt.legend()
plt.show()
\end{lstlisting}

\begin{center}
\includegraphics[width=0.7\textwidth]{figures/chap10/sin_cos.pdf}
\end{center}

\subsection{Line Styles and Colors}

\begin{lstlisting}
plt.plot(x, y, color="red")
plt.plot(x, y, color="#1f77b4")     # hex color
plt.plot(x, y, linestyle="--")       # dashed
plt.plot(x, y, linestyle=":")        # dotted
plt.plot(x, y, linewidth=2)          # thicker line

# Combined format string
plt.plot(x, y, "r--")   # red dashed
plt.plot(x, y, "b.-")   # blue dots with line
plt.plot(x, y, "go")    # green circles (no line)
\end{lstlisting}

\begin{labexample}
Plotting kinetics data:

\begin{lstlisting}
time = np.array([0, 10, 20, 30, 40, 50, 60])
concentration = np.array([1.0, 0.82, 0.67, 0.55, 0.45, 0.37, 0.30])

plt.plot(time, concentration, "ko-", label="Experimental")

# Add fitted curve
k = 0.02
t_fit = np.linspace(0, 60, 100)
c_fit = 1.0 * np.exp(-k * t_fit)
plt.plot(t_fit, c_fit, "r-", label=f"Fit: k = {k} s$^{{-1}}$")

plt.xlabel("Time (s)")
plt.ylabel("Concentration (M)")
plt.legend()
plt.savefig("kinetics.png", dpi=300)
\end{lstlisting}

\begin{center}
\includegraphics[width=0.7\textwidth]{figures/chap10/kinetics.pdf}
\end{center}
\end{labexample}

% ----------------------------------------------------------------------------
\section{Scatter Plots}
\label{sec:scatter}

\begin{lstlisting}
x = np.random.rand(50)
y = 2 * x + 0.5 + np.random.randn(50) * 0.2

plt.scatter(x, y)
plt.xlabel("Independent Variable")
plt.ylabel("Dependent Variable")
plt.show()
\end{lstlisting}

\begin{center}
\includegraphics[width=0.7\textwidth]{figures/chap10/scatter_basic.pdf}
\end{center}

\subsection{Customizing Markers}

\begin{lstlisting}
plt.scatter(x, y,
           s=50,              # marker size
           c="blue",          # color
           marker="o",        # marker shape
           alpha=0.7,         # transparency
           edgecolors="black")
\end{lstlisting}

\subsection{Color by Value}

\begin{lstlisting}
# Color points by a third variable
z = x + y
plt.scatter(x, y, c=z, cmap="viridis")
plt.colorbar(label="z value")
plt.show()
\end{lstlisting}

\begin{center}
\includegraphics[width=0.7\textwidth]{figures/chap10/scatter_colorby.pdf}
\end{center}

% ----------------------------------------------------------------------------
\section{Bar Charts}
\label{sec:bar}

\begin{lstlisting}
categories = ["A", "B", "C", "D"]
values = [23, 45, 56, 78]

plt.bar(categories, values)
plt.xlabel("Category")
plt.ylabel("Value")
plt.show()
\end{lstlisting}

\begin{center}
\includegraphics[width=0.7\textwidth]{figures/chap10/bar_basic.pdf}
\end{center}

\subsection{Grouped Bars}

\begin{lstlisting}
x = np.arange(4)
width = 0.35
treatment = [23, 45, 56, 78]
control = [20, 38, 49, 65]

plt.bar(x - width/2, treatment, width, label="Treatment")
plt.bar(x + width/2, control, width, label="Control")
plt.xticks(x, ["A", "B", "C", "D"])
plt.legend()
plt.show()
\end{lstlisting}

\begin{center}
\includegraphics[width=0.7\textwidth]{figures/chap10/bar_grouped.pdf}
\end{center}

\subsection{Error Bars}

\begin{lstlisting}
means = [23, 45, 56, 78]
stds = [2, 4, 3, 5]

plt.bar(categories, means, yerr=stds, capsize=5)
plt.ylabel("Value")
plt.show()
\end{lstlisting}

\begin{center}
\includegraphics[width=0.7\textwidth]{figures/chap10/bar_errorbars.pdf}
\end{center}

% ----------------------------------------------------------------------------
\section{Histograms}
\label{sec:histograms}

\begin{lstlisting}
data = np.random.randn(1000)  # 1000 normal random numbers

plt.hist(data, bins=30)
plt.xlabel("Value")
plt.ylabel("Frequency")
plt.show()
\end{lstlisting}

\begin{center}
\includegraphics[width=0.7\textwidth]{figures/chap10/histogram_basic.pdf}
\end{center}

\subsection{Customizing Histograms}

\begin{lstlisting}
plt.hist(data,
        bins=30,
        density=True,        # normalize to probability density
        alpha=0.7,
        color="steelblue",
        edgecolor="black")

# Overlay multiple histograms
plt.hist(data1, bins=30, alpha=0.5, label="Sample 1")
plt.hist(data2, bins=30, alpha=0.5, label="Sample 2")
plt.legend()
\end{lstlisting}

% ----------------------------------------------------------------------------
\section{Subplots}
\label{sec:subplots}

Combine multiple plots in one figure:

\begin{lstlisting}
fig, axes = plt.subplots(2, 2, figsize=(10, 8))

# axes is a 2x2 array of Axes objects
axes[0, 0].plot(x, np.sin(x))
axes[0, 0].set_title("Sine")

axes[0, 1].plot(x, np.cos(x))
axes[0, 1].set_title("Cosine")

axes[1, 0].plot(x, np.exp(-x))
axes[1, 0].set_title("Exponential Decay")

axes[1, 1].plot(x, x**2)
axes[1, 1].set_title("Quadratic")

plt.tight_layout()  # prevent overlapping
plt.savefig("subplots.png")
\end{lstlisting}

\begin{center}
\includegraphics[width=0.85\textwidth]{figures/chap10/subplots_2x2.pdf}
\end{center}

\subsection{Sharing Axes}

\begin{lstlisting}
fig, axes = plt.subplots(2, 1, sharex=True)

axes[0].plot(time, concentration)
axes[0].set_ylabel("Concentration (M)")

axes[1].plot(time, rate)
axes[1].set_ylabel("Rate (M/s)")
axes[1].set_xlabel("Time (s)")

plt.tight_layout()
\end{lstlisting}

\begin{note}
Using the \py{plt.subplots()} command is not necassary: it would also work to repeat the \py{plt.plot()} command twice and include the same array as one of the axes, but this can increase the size of the code significantly if you have more than two graphs to plot. It makes debugging easier though.
\end{note}

% ----------------------------------------------------------------------------
\section{The Object-Oriented Interface}
\label{sec:oo-interface}

For complex figures, the object-oriented interface gives more control:

\begin{lstlisting}
# Create figure and axes explicitly
fig, ax = plt.subplots(figsize=(8, 6))

ax.plot(x, y, "b-", label="Data")
ax.set_xlabel("X axis")
ax.set_ylabel("Y axis")
ax.set_title("My Plot")
ax.legend()
ax.grid(True, alpha=0.3)

fig.savefig("plot.png", dpi=300, bbox_inches="tight")
\end{lstlisting}

\begin{tip}
Use \py{plt.plot()} for quick, simple plots. Use \py{fig, ax = plt.subplots()} for anything more complex or when you need precise control.
\end{tip}

% ----------------------------------------------------------------------------
\section{Customization}
\label{sec:customization}

\subsection{Axis Limits and Ticks}

\begin{lstlisting}
ax.set_xlim(0, 10)
ax.set_ylim(-1, 1)
ax.set_xticks([0, 2, 4, 6, 8, 10])
ax.set_xticklabels(["0", "2", "4", "6", "8", "10 s"])
\end{lstlisting}

\begin{tip}
Pyplot will automatically generate ticks for your axes. If you wish to have none, use the \py{set_xticks} command with an empty set.\n
Similarly, having a \py{set_xticks{}} command for defining the values of the ticks but no \py{set_xticklabels()} command will set the labels to the values by default. Having no labels requires to explicitly provide a list of empty strings of characters using \py{set_xticklabels()}.
\end{tip}

\subsection{Log Scales}

\begin{lstlisting}
ax.set_yscale("log")
ax.set_xscale("log")

# Or use specialized plot functions
plt.semilogy(x, y)   # log y-axis
plt.semilogx(x, y)   # log x-axis
plt.loglog(x, y)     # both log
\end{lstlisting}

\subsection{Grid}

\begin{lstlisting}
ax.grid(True)
ax.grid(True, which="major", linestyle="-", alpha=0.5)
ax.grid(True, which="minor", linestyle=":", alpha=0.3)
ax.minorticks_on()
\end{lstlisting}

\subsection{Annotations}

\begin{lstlisting}
ax.annotate("Peak",
           xy=(peak_x, peak_y),          # point to annotate
           xytext=(peak_x + 1, peak_y + 0.1),  # text position
           arrowprops=dict(arrowstyle="->"))
\end{lstlisting}

% ----------------------------------------------------------------------------
\section{Style and Aesthetics}
\label{sec:style}

\subsection{Built-in Styles}

\begin{lstlisting}
print(plt.style.available)  # see all styles

plt.style.use("seaborn-v0_8-whitegrid")
# or
plt.style.use("ggplot")
\end{lstlisting}

\subsection{Custom RC Parameters}

\begin{lstlisting}
plt.rcParams["figure.figsize"] = (8, 6)
plt.rcParams["font.size"] = 12
plt.rcParams["axes.labelsize"] = 14
plt.rcParams["lines.linewidth"] = 2
\end{lstlisting}

\subsection{Publication-Ready Settings}

\begin{lstlisting}
# Settings for journal figures
plt.rcParams.update({
    "figure.figsize": (3.5, 2.5),   # single column width
    "font.size": 8,
    "axes.labelsize": 9,
    "axes.titlesize": 9,
    "legend.fontsize": 7,
    "xtick.labelsize": 7,
    "ytick.labelsize": 7,
    "font.family": "sans-serif",
    "figure.dpi": 300
})
\end{lstlisting}

% ----------------------------------------------------------------------------
\section{Specialized Plots}
\label{sec:specialized}

\subsection{Contour Plots}

\begin{lstlisting}
x = np.linspace(-3, 3, 100)
y = np.linspace(-3, 3, 100)
X, Y = np.meshgrid(x, y)
Z = np.sin(X) * np.cos(Y)

plt.contour(X, Y, Z, levels=10)
plt.colorbar()
\end{lstlisting}

\begin{center}
\includegraphics[width=0.6\textwidth]{figures/chap10/contour.pdf}
\end{center}

Filled contours use color to show values between contour lines:

\begin{lstlisting}
# Filled contours
plt.contourf(X, Y, Z, levels=20, cmap="RdBu_r")
plt.colorbar()
\end{lstlisting}

\begin{center}
\includegraphics[width=0.6\textwidth]{figures/chap10/contourf.pdf}
\end{center}

\begin{note}
This colour scale feels intuitive especially for temperatures, but it yields a poor contrast if printed, especially in black & white, and is not really colorblind-friendly. That's why "viridis", a scale from purple to yellow, is preferred and used in the other examples of this section.
\end{note}

\subsection{Heatmaps}

\begin{lstlisting}
data = np.random.rand(10, 10)

plt.imshow(data, cmap="viridis", aspect="auto")
plt.colorbar(label="Value")
plt.xlabel("Column")
plt.ylabel("Row")
\end{lstlisting}

\begin{center}
\includegraphics[width=0.6\textwidth]{figures/chap10/heatmap.pdf}
\end{center}

\subsection{Box Plots}

\begin{lstlisting}
data = [np.random.randn(100) + i for i in range(4)]

plt.boxplot(data, labels=["A", "B", "C", "D"])
plt.ylabel("Value")
\end{lstlisting}

\begin{center}
\includegraphics[width=0.6\textwidth]{figures/chap10/boxplot.pdf}
\end{center}

% ----------------------------------------------------------------------------
\section{A Complete Figure Example}
\label{sec:complete-figure}

\begin{lstlisting}
import matplotlib.pyplot as plt
import numpy as np

# Generate sample spectroscopy data
wavelengths = np.linspace(400, 700, 150)
spectrum1 = np.exp(-((wavelengths - 500)**2) / 1000) * 0.8
spectrum2 = np.exp(-((wavelengths - 550)**2) / 800) * 0.6
spectrum3 = np.exp(-((wavelengths - 480)**2) / 1200) * 0.9

# Create figure
fig, axes = plt.subplots(2, 1, figsize=(8, 6), sharex=True)

# Top panel: raw spectra
ax1 = axes[0]
ax1.plot(wavelengths, spectrum1, label="Sample A")
ax1.plot(wavelengths, spectrum2, label="Sample B")
ax1.plot(wavelengths, spectrum3, label="Sample C")
ax1.set_ylabel("Absorbance (AU)")
ax1.legend(loc="upper right")
ax1.set_title("UV-Vis Spectra Comparison")

# Bottom panel: difference spectrum
ax2 = axes[1]
difference = spectrum1 - spectrum2
ax2.plot(wavelengths, difference, "k-")
ax2.axhline(y=0, color="gray", linestyle="--", alpha=0.5)
ax2.set_ylabel("Difference (A - B)")
ax2.set_xlabel("Wavelength (nm)")

# Add peak annotations
peak_idx = np.argmax(spectrum1)
ax1.annotate(f"Peak: {wavelengths[peak_idx]:.0f} nm",
            xy=(wavelengths[peak_idx], spectrum1[peak_idx]),
            xytext=(wavelengths[peak_idx] + 50, spectrum1[peak_idx] + 0.05),
            arrowprops=dict(arrowstyle="->", color="gray"))

plt.tight_layout()
plt.savefig("spectra_comparison.pdf", bbox_inches="tight")
plt.show()
\end{lstlisting}

\begin{center}
\includegraphics[width=0.85\textwidth]{figures/chap10/spectra_complete.pdf}
\end{center}

% ----------------------------------------------------------------------------
\section{Chapter Summary}
\label{sec:matplotlib-summary}

Matplotlib creates publication-quality figures:

\begin{itemize}
    \item \textbf{Basic:} \py{plt.plot()}, \py{plt.scatter()}, \py{plt.bar()}, \py{plt.hist()}
    \item \textbf{Labels:} \py{plt.xlabel()}, \py{plt.ylabel()}, \py{plt.title()}, \py{plt.legend()}
    \item \textbf{Subplots:} \py{fig, axes = plt.subplots(rows, cols)}
    \item \textbf{Customization:} colors, line styles, markers, axes limits
    \item \textbf{Saving:} \py{plt.savefig("file.pdf", dpi=300)}
    \item \textbf{Styles:} \py{plt.style.use()}, \py{plt.rcParams}
\end{itemize}

\begin{ponder}
\begin{enumerate}
    \item When would you use \py{plt.plot()} vs \py{plt.scatter()}?
    \item Why save figures as PDF rather than PNG for publications?
    \item How would you create a figure with 3 subplots in a row?
    \item What's the advantage of the object-oriented interface (\py{fig, ax = plt.subplots()})?
\end{enumerate}
\end{ponder}
